\subsection{模型总结}

三问采用统一的单位约定、Snell折射关系和带符号Fresnel界面系数，形成“正演模型--参数反演--多光束影响分析”的递进结构。问题一建立Sellmeier--Drude--Fresnel双光束正演模型；问题二在$2500$--$3300\,\mathrm{cm^{-1}}$弱吸收波段内对两个角度独立反演并取等权平均，得到SiC外延层厚度$7.7398\,\mu\mathrm m$；问题三引入双振子--Drude复折射率与Airy多光束模型，得到硅外延层厚度$3.2175\,\mu\mathrm m$，并通过同参双光束截断比较、模型误差厚度传播、Jacobian--SVD及SiC高阶反射量级分析多光束影响。综合上述结果，SiC弱吸收波段内的高阶反射不足以引起可辨的厚度修正。

\subsection{模型优点}

\begin{enumerate}[label=\textbf{\arabic*.}]
\item \textbf{物理结构前后一致。} 问题一至问题三均从分层介质传播和Fresnel界面条件出发，模型复杂度随题目物理条件逐步增加：SiC弱吸收区使用实折射率近似，硅明显吸收区恢复复折射率与传播衰减。带符号反射系数直接包含界面相位反转，避免与传播相位重复计入固定半波损失。

\item \textbf{厚度反演降低了初值与附属参数的影响。} 问题二采用"低维搜索--候选去重--完整精修"两阶段策略，保留干涉级次诱导的多个局部极小，并以残差竞争确定各角度厚度。最优盆地已在固定Sellmeier阶段形成，释放七个参数后仍保持最低RMSE，说明厚度级次不依赖额外自由度；两角度独立反演结果相近，厚度尺度具有较好一致性。

\item \textbf{参数约束具有可辨识性依据。} 问题三在最优Airy解处构造尺度化Jacobian并进行SVD，区分厚度稳定性与附属参数的弱可辨识性；释放背景参数后条件数显著恶化，说明固定双振子背景和有效质量可减少有限波段内的近共线自由度。

\item \textbf{多光束影响被量化到厚度尺度。} 问题三在Airy最优参数下计算双光束截断谱形差，并将$\Delta R$通过截断模型Jacobian传播到厚度方向；两角度局部厚度偏移均不超过$1.2\times10^{-3}\,\mu\mathrm m$。代入问题二SiC参数后，第三束相对表面反射束的最大强度比仅约$4.85\times10^{-7}\%$，与弱吸收区采用双光束近似的判断一致。
\end{enumerate}

\subsection{模型局限}

\begin{enumerate}[label=\textbf{\arabic*.}]
\item \textbf{观测维度仍可扩展。} 本文利用两个入射角光谱完成厚度反演，并以跨角度结果检验厚度尺度的一致性；若后续增加入射角、偏振或重复测量，可进一步丰富交叉验证信息。

\item \textbf{材料与界面先验可进一步细化。} 受附件信息限制，部分光源与界面条件采用有效参数处理；结合独立折射率、相干长度或粗糙度测量后，可继续完善谱形细节的物理刻画。
\end{enumerate}

\subsection{改进方向}

若能获取更多入射角、偏振分辨光谱或独立折射率测量数据，可构造带实验先验的多角度联合约束，在保持参数可辨识性的前提下进一步降低厚度反演的不确定性；若取得晶圆多位置测量数据，则可由单点厚度反演扩展为二维厚度分布估计，分别报告空间非均匀性与光谱拟合误差。对于多光束模型，还可在获得光源相干长度和界面粗糙度信息后引入有限相干与散射衰减项，进一步提升吸收波段的物理描述精度。
